\documentclass[10pt,a4paper]{amsart}
\usepackage{amsmath,amssymb,amsthm,graphicx}
\usepackage{verbatim}
\newenvironment{reference}{\comment}{\endcomment}
%\newenvironment{reference}{}{}
\newenvironment{slogan}{\comment}{\endcomment}
\newenvironment{history}{\comment}{\endcomment}

% For commutative diagrams we use Xy-pic
\usepackage[all]{xy}

% We use 2cell for 2-commutative diagrams.
\xyoption{2cell}
% Theorem environments.
%
\theoremstyle{plain}
\newtheorem{theorem}[subsection]{Theorem}
\newtheorem{proposition}[subsection]{Proposition}
\newtheorem{lemma}[subsection]{Lemma}

\theoremstyle{definition}
\newtheorem{definition}[subsection]{Definition}
\newtheorem{example}[subsection]{Example}
\newtheorem{exercise}[subsection]{Exercise}
\newtheorem{situation}[subsection]{Situation}

\theoremstyle{remark}
\newtheorem{remark}[subsection]{Remark}
\newtheorem{remarks}[subsection]{Remarks}

\numberwithin{equation}{subsection}

% Macros
%
\def\lim{\mathop{\mathrm{lim}}\nolimits}
\def\colim{\mathop{\mathrm{colim}}\nolimits}
\def\Spec{\mathop{\mathrm{Spec}}}
\def\Hom{\mathop{\mathrm{Hom}}\nolimits}
\def\Ext{\mathop{\mathrm{Ext}}\nolimits}
\def\SheafHom{\mathop{\mathcal{H}\!\mathit{om}}\nolimits}
\def\SheafExt{\mathop{\mathcal{E}\!\mathit{xt}}\nolimits}
\def\Sch{\mathit{Sch}}
\def\Mor{\mathop{\mathrm{Mor}}\nolimits}
\def\Ob{\mathop{\mathrm{Ob}}\nolimits}
\def\Sh{\mathop{\mathit{Sh}}\nolimits}
\def\NL{\mathop{N\!L}\nolimits}
\def\CH{\mathop{\mathrm{CH}}\nolimits}
\def\proetale{{pro\text{-}\acute{e}tale}}
\def\etale{{\acute{e}tale}}
\def\QCoh{\mathit{QCoh}}
\def\Ker{\mathop{\mathrm{Ker}}}
\def\Im{\mathop{\mathrm{Im}}}
\def\Coker{\mathop{\mathrm{Coker}}}
\def\Coim{\mathop{\mathrm{Coim}}}

% Boxtimes
%
\DeclareMathSymbol{\boxtimes}{\mathbin}{AMSa}{"02}

%
% Macros for moduli stacks/spaces
%
\def\QCohstack{\mathcal{QC}\!\mathit{oh}}
\def\Cohstack{\mathcal{C}\!\mathit{oh}}
\def\Spacesstack{\mathcal{S}\!\mathit{paces}}
\def\Quotfunctor{\mathrm{Quot}}
\def\Hilbfunctor{\mathrm{Hilb}}
\def\Curvesstack{\mathcal{C}\!\mathit{urves}}
\def\Polarizedstack{\mathcal{P}\!\mathit{olarized}}
\def\Complexesstack{\mathcal{C}\!\mathit{omplexes}}
% \Pic is the operator that assigns to X its picard group, usage \Pic(X)
% \Picardstack_{X/B} denotes the Picard stack of X over B
% \Picardfunctor_{X/B} denotes the Picard functor of X over B
\def\Pic{\mathop{\mathrm{Pic}}\nolimits}
\def\Picardstack{\mathcal{P}\!\mathit{ic}}
\def\Picardfunctor{\mathrm{Pic}}
\def\Deformationcategory{\mathcal{D}\!\mathit{ef}}
\usepackage{cleveref}
\crefname{originallemma}{original lemma}{original lemmas}
\Crefname{originallemma}{Original lemma}{Original lemmas}
\makeatletter
\newlabel{lemma-quotient-stack-functorial}{{046Q}{source}{}{sourcepointer-1}{}}
\newlabel{lemma-quotient-stack-functorial@cref}{{[originallemma][0][]046Q}{[1][0][]source}{}{}{}}
\makeatother
\crefname{originallemma}{original lemma}{original lemmas}
\Crefname{originallemma}{Original lemma}{Original lemmas}
\makeatletter
\newlabel{lemma-cartesian-square-of-morphism}{{04Y4}{source}{}{sourcepointer-2}{}}
\newlabel{lemma-cartesian-square-of-morphism@cref}{{[originallemma][0][]04Y4}{[1][0][]source}{}{}{}}
\makeatother
\crefname{originallemma}{original lemma}{original lemmas}
\Crefname{originallemma}{Original lemma}{Original lemmas}
\makeatletter
\newlabel{lemma-quotient-stack-restrict-equivalence}{{046T}{source}{}{sourcepointer-3}{}}
\newlabel{lemma-quotient-stack-restrict-equivalence@cref}{{[originallemma][0][]046T}{[1][0][]source}{}{}{}}
\makeatother
\title{Cartesian R-square yields a 2-fibre product square of quotient stacks}
\author{The Stacks Project Authors}
\date{stacks-project:spaces-groupoids.tex:local-snapshot-20261012; GNU FDL 1.2 or later}
\begin{document}
\maketitle
\noindent\textit{Editorial note.} All delivered mathematical text is the author's own, copied verbatim from the pinned Stacks Project chapter file spaces-groupoids.tex; no statement or proof step was rewritten, reordered or removed.\par\medskip
\noindent\textit{Editorial note.} Span recomputation by label pairing (LANE\_BRIEF 10.2): in this chapter the author's proof environment follows the statement environment as a standalone block; the spec covers the complete environment pair (statement span, proof span) as recorded in the pool candidate.\par\medskip
\ifdefined\hypertarget\hypertarget{sourcepointer-1}{}\fi
\noindent\textit{Source reference.} Original lemma 046Q: sources/stacks-chow/tags.txt, source lines 5481--5481. Cross-chapter standard prerequisite cited by the author inside the counted span; it is used as a pointer and its own proof is not part of this item.\par\medskip
\ifdefined\hypertarget\hypertarget{sourcepointer-2}{}\fi
\noindent\textit{Source reference.} Original lemma 04Y4: sources/stacks-chow/tags.txt, source lines 6417--6417. Cross-chapter standard prerequisite cited by the author inside the counted span; it is used as a pointer and its own proof is not part of this item.\par\medskip
\ifdefined\hypertarget\hypertarget{sourcepointer-3}{}\fi
\noindent\textit{Source reference.} Original lemma 046T: sources/stacks-chow/tags.txt, source lines 5484--5484. Cross-chapter standard prerequisite cited by the author inside the counted span; it is used as a pointer and its own proof is not part of this item.\par\medskip
\begin{lemma}
\label{lemma-criterion-fibre-product}
Notation and assumption as in
Lemma \ref{lemma-quotient-stack-functorial}.
Assume that
$$
\xymatrix{
R \ar[d]_s \ar[r]_f & R' \ar[d]^{s'} \\
U \ar[r]^f & U'
}
$$
is cartesian. Then
$$
\xymatrix{
\mathcal{S}_U \ar[d] \ar[r] & [U/R] \ar[d]^{[f]} \\
\mathcal{S}_{U'} \ar[r] & [U'/R']
}
$$
is a $2$-fibre product square.
\end{lemma}

\begin{proof}
Applying the inverse isomorphisms $i : R \to R$ and $i' : R' \to R'$
to the (first) cartesian diagram of the statement of the lemma we see
that
$$
\xymatrix{
R \ar[d]_t \ar[r]_f & R' \ar[d]^{t'} \\
U \ar[r]^f & U'
}
$$
is cartesian as well. By
Lemma \ref{lemma-cartesian-square-of-morphism}
we have a $2$-fibre square
$$
\xymatrix{
[U''/R''] \ar[d] \ar[r] & [U/R] \ar[d] \\
\mathcal{S}_{U'} \ar[r] & [U'/R']
}
$$
where $U'' = U \times_{f, U', t'} R'$ and
$R'' = R \times_{f \circ s, U', t'} R'$. By the above
we see that $(t, f) : R \to U''$ is an isomorphism, and that
$$
R'' =
R \times_{f \circ s, U', t'} R' =
R \times_{s, U} U \times_{f, U', t'} R' =
R \times_{s, U, t} \times R.
$$
Explicitly the isomorphism $R \times_{s, U, t} R \to R''$ is given by
the rule $(r_0, r_1) \mapsto (r_0, f(r_1))$.
Moreover, $s'', t'', c''$ translate into the maps
$$
R \times_{s, U, t} R \to R,
\quad
s''(r_0, r_1) = r_1, \quad t''(r_0, r_1) = c(r_0, r_1)
$$
and
$$
\begin{matrix}
c'' : &
(R \times_{s, U, t} R) \times_{s'', R, t''} (R \times_{s, U, t} R)
&
\longrightarrow
&
R \times_{s, U, t} R, \\
&
((r_0, r_1), (r_2, r_3)) &
\longmapsto & (c(r_0, r_2), r_3).
\end{matrix}
$$
Precomposing with the isomorphism
$$
R \times_{s, U, s} R \longrightarrow R \times_{s, U, t} R,
\quad
(r_0, r_1) \longmapsto (c(r_0, i(r_1)), r_1)
$$
we see that $t''$ and $s''$ turn into
$\text{pr}_0$ and $\text{pr}_1$ and that $c''$ turns into
$\text{pr}_{02} :
R \times_{s, U, s} R \times_{s, U, s} R \to R \times_{s, U, s} R$.
Hence we see that there is an isomorphism
$[U''/R''] \cong [R/R \times_{s, U, s} R]$ where as a groupoid
in algebraic spaces $(R, R \times_{s, U, s} R, s'', t'', c'')$
is the restriction of the trivial groupoid
$(U, U, \text{id}, \text{id}, \text{id})$ via $s : R \to U$.
Since $s : R \to U$ is a surjection of fppf sheaves (as it has a
right inverse) the morphism
$$
[U''/R''] \cong [R/R \times_{s, U, s} R]
\longrightarrow
[U/U] = \mathcal{S}_U
$$
is an equivalence by
Lemma \ref{lemma-quotient-stack-restrict-equivalence}.
This proves the lemma.
\end{proof}

\end{document}
